% PDF 124--127；表 3.1 见 PDF 128--129。
\subsection{基于提示的知识提取}
语言模型在预训练中既学习一般性的语义知识，也学到无数具体的事实性陈述，这些知识通常称为参数化知识（parametric knowledge）。图~\ref{fig:ch03-formal-prompt} 的提示任务只需了解词语之间的语义关系即可完成。再考虑另一个问题，“What city was Dante Alighieri born in?”，也就是问但丁·阿利吉耶里（Dante Alighieri）出生在哪座城市。仅凭语义无法回答这个问题，模型必须接触过具体的事实信息。大语言模型预训练语料库的巨大规模使这一点成为可能。例如，如果语料库含有但丁的维基百科词条，模型权重就会受到这样的句子影响。“但丁出生在佛罗伦萨共和国的佛罗伦萨，那里现在位于意大利。”恰当设计的提示可以让模型复述这类参数化知识。例如，
\begin{equation}
\text{\ttfamily Dante was born in [MASK]}
\tag{3.4}
\end{equation}
由于模型训练时接触过能回答这个问题的数据，就会认为“Florence”很可能填在掩码词元的位置。这说明了大语言模型基于提示的推理的一项基本潜力，即可以将大语言模型作为知识库（knowledge base）使用。

构建传统知识库需要从不同数据源提取关系知识，也需要开发实体抽取、共指消解（coreference resolution）、实体链接（entity linking）和关系抽取（relation extraction）等自然语言处理流程（\refcite{ch03-petroni}{Petroni et al., 2019）}。各环节都有难点，错误难以避免。错误沿流程传播和积累，可能显著影响所建知识库的实用性（\refcite{ch03-petroni}{Petroni et al., 2019）}。图~\ref{fig:ch03-kb} 从概念上比较了两种做法。

\begin{figure}[!htbp]
\centering
\scalebox{.99}{% 重绘自原书第 95 页的图 3.9。
\begin{tikzpicture}[
  x=.97cm,font=\footnotesize\sffamily,
  head/.style={font=\footnotesize\sffamily\bfseries},
  vertex/.style={circle,inner sep=0pt,minimum size=3.2mm,draw=black!76,line width=.7pt},
  wire/.style={draw=black!68,line width=.55pt},
  flow/.style={wire,-{Stealth[length=1.4mm]}},
  memory/.style={fill=yellow!22,align=center,minimum width=23mm,minimum height=8mm,inner sep=2pt}]
\node[head] at (1.95,4.96) {记忆};
\node[head] at (5.97,4.96) {查询};
\node[head] at (9.80,4.96) {答案};
\node[head,anchor=east] at (.05,3.47) {知识库};
\node[vertex,fill=blue!35] (base) at (.40,3.47) {};
\node[vertex,fill=green!30] (top) at (.95,4.12) {};
\node[vertex,fill=green!30] (dante) at (1.85,3.47) {};
\node[vertex,fill=green!30] (bottom) at (1.19,2.83) {};
\node[vertex,fill=yellow!63] (topentity) at (2.65,4.20) {};
\node[vertex,fill=yellow!63] (florence) at (2.65,2.66) {};
\draw[wire] (base)--(top)--(bottom)--(base)--(dante)--(topentity);
\draw[wire] (dante) -- node[font=\tiny\ttfamily,fill=white,inner sep=.8pt,sloped] {born-in} (florence);
\node[anchor=west,font=\scriptsize] at (2.04,3.49) {Dante};
\node[anchor=west,font=\scriptsize] at (2.83,2.65) {Florence};

\node[memory] (kb) at (5.97,3.41) {访问符号记忆};
\node[font=\scriptsize\ttfamily] (rel) at (5.97,4.23) {(Dante, born-in, X)};
\node at (9.80,3.41) {Florence};
\draw[wire] (3.53,3.41) -- (kb.west);
\draw[flow] (rel.south) -- (kb.north);
\draw[flow] (kb.east) -- (9.11,3.41);

\draw[draw=black!62,densely dotted,line width=1pt] (-.62,2.13) -- (10.48,2.13);
\node[head,anchor=east] at (.05,.88) {语言模型};
\foreach \offset in {.32,.16,0} {
  \begin{scope}[shift={(\offset,\offset)}]
  \filldraw[draw=black!25,fill=black!6] (.58,.35) rectangle (1.82,1.59);
  \foreach \x/\y/\tone in {.82/.59/blue,.82/.97/green,.82/1.35/green,
   1.20/.59/green,1.20/.97/yellow,1.20/1.35/blue,
   1.58/.59/green,1.58/.97/yellow,1.58/1.35/yellow} {
    \filldraw[draw=black!67,fill=\tone!28] (\x,\y) circle[radius=1.88mm];
  }
  \end{scope}
}
\node[font=\scriptsize] at (1.28,.03) {如 ELMo/BERT};
\node[memory] (lm) at (5.97,.80) {访问神经语言\\模型记忆};
\node[font=\scriptsize\ttfamily] (prompt) at (5.97,1.68) {"Dante was born in [MASK]"};
\node at (9.80,.80) {Florence};
\draw[wire] (3.53,.80) -- (lm.west);
\foreach \x/\y in {4.96/5.28,5.55/5.69,6.11/6.11,6.57/6.46,7.08/6.83} {
  \draw[flow] (\x,1.49) -- (\y,1.21);
}
\draw[flow] (lm.east) -- (9.11,.80);
\end{tikzpicture}
}
\caption{向传统知识库和语言模型查询事实信息。虚线上方是为查询实体关系而专门设计的知识库，本例关系是“但丁出生在佛罗伦萨”。下方的语言模型原本用于根据上下文预测被遮盖的词。因此，可借助提示让它报告预训练时接触过的事实。}
\label{fig:ch03-kb}
\end{figure}

\begin{keypoint}
与使用传统知识库相比，给大语言模型提供提示有几个主要优点。首先，它不需要预先定义模式（schema-free），关系知识是预训练的涌现结果，存于语言模型内部，无须专为这项任务开发模型。其次，现代预训练语料覆盖的信息范围很广，模型的泛化能力也很强。理论上，同一个语言模型可以支持多种知识应用和范围更广的常见自然语言处理任务。
\end{keypoint}

\begin{practicaltip}
基于大语言模型提取知识也面临很大的挑战和风险。要解释模型回答并不容易，因为往往无法追查是哪些因果过程使基于提示的推理产生某个特定回答。回答的准确性也需要验证。我们试图引出的知识，是预训练过程中涌现的，模型并未接受专门训练以学习这些事实。预训练数据的规模也大到难以验证知识或按质量排序。若数据取自互联网，同样缺乏对知识来源和可靠性的审查，可能使模型学到不真实的“事实”。因此，用户应保持合理的怀疑态度，采用适当的评估方法防范这些错误。最后，基于提示的知识库能否给出一致结果，很大程度上取决于提示设计的质量。第 3.3 节将详细讨论提示优化的方法。
\end{practicaltip}

\subsection{基于提示的学习在各类自然语言处理任务中的应用}
许多自然语言处理任务都适合用提示来完成。本节列出常见任务，并介绍它们的输入、模板、提示及答案映射。理解这些要素以及它们之间的关系，有助于了解如何把各项自然语言处理任务有效转化成基于提示的学习问题并加以求解。

我们将这些任务分成三类。
\begin{enumerate}
\item \textbf{文本分类。}为给定的输入句子分配适当的类别标签。提示接收输入句子，留出专门位置供模型生成中间答案，随后可把中间答案映射为分类标签。
\item \textbf{标注（tagging）。}为给定文本中的个别元素，如词或短语，分配标签或标记。提示先包含待标注元素所在的文本，再就该元素具体提问，给出若干选项供模型选择。
\item \textbf{文本生成。}根据提示中的任务生成一段文本，通常包含不止一个词元。提示给出相关上下文，如待摘要的段落或待翻译的句子，再用明确指令告诉模型如何处理这段上下文。
\end{enumerate}

表~\ref{tab:ch03-tasks} 列出三类共 17 项任务，简述各项任务，并给出可以用于完成该任务的示例输入、模板和答案空间。多样的用例体现了基于提示的学习的灵活性。不过，提示需要根据每项任务精心设计。下一节将把提示设计拆分为几个可以分别优化的方面，以改善基于提示的学习的表现。

% 按原书第 98--99 页重排表 3.1。
\begingroup
\small
\setlength{\tabcolsep}{4pt}
\setlength{\LTcapwidth}{\linewidth}
\renewcommand{\arraystretch}{1.07}
\newcommand{\taskexample}[3]{%
  \textbf{输入}\enspace #1\par
  \textbf{模板}\enspace #2\par
  \textbf{答案}\enspace #3}
\newcommand{\inputslots}{[$\symbf{x}_1,\ldots,\symbf{x}_n$]}
\begin{longtable}{@{}>{\raggedright\arraybackslash}p{.37\linewidth}>{\raggedright\arraybackslash}p{.60\linewidth}@{}}
\caption{基于提示的自然语言处理方法。每行左侧给出任务及定义，右侧给出示例，其中包括待处理的输入句子、用于基于提示的推理的模板，以及一种可能的答案空间。任务分为文本分类、标注与文本生成三类。}
\label{tab:ch03-tasks}\\
\endfirsthead
\multicolumn{2}{@{}l}{\textbf{表~\ref{tab:ch03-tasks}（续）}}\\
\toprule
\textbf{任务} & \textbf{示例}\\
\midrule
\endhead
\midrule
\multicolumn{2}{r@{}}{\textit{续下页}}\\
\endfoot
\bottomrule
\endlastfoot
\multicolumn{2}{@{}l}{\textbf{文本分类}}\\
\toprule
\textbf{任务} & \textbf{示例}\\
\midrule
\textbf{情感分析}\enspace 将文本的情感判为正面、负面或中性。 &
\taskexample{I hate this movie.}{\texttt{[x]} It was a \texttt{[z]} movie.}{great, terrible, \ldots}\\
\midrule
\textbf{作者归属识别}\enspace 从预先给定的作者集合中判断给定文本的作者。 &
\taskexample{It was the best of times, it was the worst of times, it was the age of wisdom, it was the age of foolishness}
{The author of \texttt{[x]} is most likely \texttt{[z]}.}{Dickens, Carroll, Austin, \ldots}\\
\midrule
\textbf{垃圾消息检测}\enspace 判断邮件或短信是否属于垃圾消息。 &
\taskexample{Congratulations! You have won! Click here to claim your free vacation.}
{This message: \texttt{[x]} is classified as \texttt{[z]}.}{Spam, Non-Spam}\\
\midrule
\textbf{情绪分类}\enspace 从预定义的情绪集合中，判断文本表达了哪种情绪。 &
\taskexample{I just won the lottery!}{This text: \texttt{[x]} expresses the emotion \texttt{[z]}.}
{anger, surprise, sadness, happiness}\\
\midrule
\textbf{意图检测}\enspace 判断用户发出查询或消息的意图，常用于聊天机器人和虚拟助手。 &
\taskexample{What’s the weather like today?}{\texttt{[x]} The user’s intent is \texttt{[z]}.}
{\texttt{get\_weather}, \texttt{set\_alarm}}\\
\midrule
\textbf{语种识别}\enspace 判断给定文本以何种语言写成。 &
\taskexample{¿Cómo estás?}{\texttt{[x]} The language is \texttt{[z]}.}{Spanish, French, \ldots}\\
\midrule
\textbf{仇恨言论检测}\enspace 判断给定文本是否含有仇恨言论。 &
\taskexample{I can’t stand them.}{\texttt{[x]} The text contains \texttt{[z]} speech.}{hate, non-hate}\\
\addlinespace[5pt]
\multicolumn{2}{@{}l}{\textbf{标注}}\\
\toprule
\textbf{任务} & \textbf{示例}\\
\midrule
\textbf{词性标注（POS）}\enspace 给词分配语法类别，如名词、动词、形容词和副词。 &
\taskexample{She is running in the park.}{In the sentence \inputslots, the word [$\symbf{x}_i$] has POS-tag [$\symbf{z}_j$].}{noun, verb, adjective, \ldots}\\
\midrule
\textbf{命名实体识别（NER）}\enspace 识别并分类文本中提到的人物、日期、地点、机构等实体。 &
\taskexample{John met Mary in London.}{In the sentence \inputslots, the word [$\symbf{x}_i$] the named entity label is [$\symbf{z}_j$].}{location, organization, \ldots}\\
\midrule
\textbf{组块分析或浅层句法分析}\enspace 依据语法结构把相邻的词或词元组成较大的单元，即组块，例如名词短语或动词短语。 &
\taskexample{She is running in the park.}{In the sentence \inputslots, the word [$\symbf{x}_i$] the chunk label is [$\symbf{z}_j$].}{’B-VP’（动词短语起点），’I-VP’（动词短语内部），\ldots}\\
\midrule
\textbf{依存句法分析}\enspace 找出句子中词语之间的句法依存关系，包括标注主语、宾语、修饰语等，并显示它们之间的关系。 &
\taskexample{She is running in the park.}{In the sentence \inputslots, the word [$\symbf{x}_i$] the dependency relation is [$\symbf{z}_j$].}{’nsubj’（名词性主语），’root’（句子的根），’dobj’（直接宾语），\ldots}\\
\midrule
\textbf{成分句法分析或短语结构分析}\enspace 识别句子的成分结构，即词如何组成名词短语（NP）、动词短语（VP）等语法短语。 &
\taskexample{She is running in the park.}{In the sentence \inputslots, the word [$\symbf{x}_i$] the constituent category is [$\symbf{z}_j$].}{’NP’（名词短语），’PP’（介词短语），’VP’（动词短语），\ldots}\\
\midrule
\textbf{语义角色标注（SRL）}\enspace 按照词或短语与谓词（通常为动词）之间的语义关系，为它们分配施事、工具等角色。 &
\taskexample{John gave Mary a book.}{In the sentence \inputslots, the word [$\symbf{x}_i$] has semantic role [$\symbf{z}_j$].}{Agent, Theme, Location, \ldots}\\
\midrule
\textbf{共指消解}\enspace 识别文本中指向同一实体的词语或短语，并将它们关联起来。 &
\taskexample{Jane is a talented software engineer. She was recently promoted to team lead.}{In the text with words: \inputslots, does the word [$\symbf{x}_i$] refer to the word $\symbf{x}_j$?}{Yes, No}\\
\addlinespace[5pt]
\multicolumn{2}{@{}l}{\textbf{文本生成}}\\
\toprule
\textbf{任务} & \textbf{示例}\\
\midrule
\textbf{摘要生成}\enspace 给定一篇长文本，生成保留原文要点或关键信息的较短版本。 &
\taskexample{<Long text to be summarized.>}{Please provide a summary for the following text: \texttt{[x]}. Summary: \texttt{[z]}.}{<summarized version of the long text>}\\
\midrule
\textbf{问答}\enspace 给定问题和上下文，依据上下文中的信息生成答案。 &
\taskexample{<Context or passage>, <question>}{Here is the context: \texttt{[x]} What is the answer to the question: \texttt{[w]} ? Answer: \texttt{[z]} .}{<answer to the question based on the context>}\\
\midrule
\textbf{机器翻译}\enspace 将文本从一种语言译成另一种语言，保留原意和语境。 &
\taskexample{¿Cómo estás?}{Translate the following text from the source language to the target language: \texttt{[x]} Translation: \texttt{[z]}.}{<translated text in target language>}\\
\end{longtable}
\endgroup

